% !TeX program = lualatex
\documentclass[12pt,aspectratio=169]{beamer}
\usepackage[utf8]{inputenc}
\usepackage[T1]{fontenc}
\usepackage{amsmath,amsfonts,amssymb}
\usepackage{graphicx}
\usepackage{xcolor}
\usepackage{hyperref}
\usepackage{anyfontsize}
\usepackage{lmodern}

% ====== EXTREME FONT REDUCTION ======
\renewcommand{\tiny}{\fontsize{4}{5}\selectfont}
\renewcommand{\scriptsize}{\fontsize{5}{6}\selectfont}
\renewcommand{\footnotesize}{\fontsize{6}{7}\selectfont}
\renewcommand{\small}{\fontsize{7}{8}\selectfont}
\renewcommand{\normalsize}{\fontsize{8}{9}\selectfont}
\renewcommand{\large}{\fontsize{9}{10}\selectfont}
\renewcommand{\Large}{\fontsize{10}{11}\selectfont}
\renewcommand{\LARGE}{\fontsize{11}{13}\selectfont}
\renewcommand{\huge}{\fontsize{12}{14}\selectfont}
\renewcommand{\Huge}{\fontsize{14}{17}\selectfont}

% ====== COLOR THEME ======
\definecolor{redcorrect}{RGB}{200,30,30}
\definecolor{bluecorrect}{RGB}{30,80,200}
\definecolor{greencheck}{RGB}{0,150,50}
\definecolor{lightgray}{RGB}{245,245,245}
\definecolor{darkgray}{RGB}{50,50,50}
\definecolor{softyellow}{RGB}{255,248,200}
\definecolor{steppurple}{RGB}{150,50,200}
\definecolor{deepblue}{RGB}{20,60,150}
\definecolor{darkred}{RGB}{150,20,20}
\definecolor{orangewarning}{RGB}{255,140,0}

\setbeamercolor{title}{fg=white,bg=redcorrect}
\setbeamercolor{frametitle}{fg=white,bg=redcorrect}
\setbeamercolor{block title}{fg=white,bg=bluecorrect}
\setbeamercolor{block body}{fg=darkgray,bg=lightgray}
\setbeamercolor{normal text}{fg=darkgray}
\usecolortheme{default}

% ====== CUSTOM COMMANDS ======
\newcommand{\correct}[1]{\textcolor{greencheck}{\textbf{\checkmark}} #1}
\newcommand{\keypoint}[1]{\textcolor{deepblue}{\textbf{#1}}}
\newcommand{\hardconcept}[1]{\textcolor{darkred}{\textbf{#1}}}
\newcommand{\tipbox}[1]{\fcolorbox{bluecorrect}{softyellow}{\parbox{0.88\textwidth}{\centering #1}}}
\newcommand{\stepbox}[1]{%
	\begin{center}
		\fcolorbox{darkgray}{white}{%
			\parbox{0.90\textwidth}{\vspace{0.05cm} #1 \vspace{0.05cm}}%
		}
	\end{center}
}
\newcommand{\wrongbox}[1]{\fcolorbox{redcorrect}{red!10}{\parbox{0.88\textwidth}{\centering \textcolor{redcorrect}{\textbf{\times}} #1}}}
\newcommand{\highlight}[1]{\textcolor{redcorrect}{\textbf{#1}}}
\newcommand{\bluetext}[1]{\textcolor{bluecorrect}{\textbf{#1}}}

\begin{document}
	
	% ================================================================
	% SLIDE 1: TITLE
	% ================================================================
	\begin{frame}
		\centering
		\vspace{0.8cm}
		{\fontsize{16}{19}\selectfont \textbf{Building an Anti-Financial Crime Compliance Program}}\\[0.15cm]
		{\fontsize{12}{14}\selectfont Domain C — Complete Q\&A Review}\\[0.3cm]
		{\fontsize{9}{11}\selectfont Understanding the \textit{Regulatory Principles} Behind Each Answer}\\[0.3cm]
		{\fontsize{8}{10}\selectfont Compliance Training Department}
		\vspace{0.2cm}
		\rule{4cm}{0.5pt}
	\end{frame}
	
	% ================================================================
	% SLIDE 2: AGENDA
	% ================================================================
	\begin{frame}
		\frametitle{\fontsize{11}{13}\selectfont Agenda — Domain C Topics}
		\scriptsize
		\begin{enumerate}
			\item Effective AFC Program and Core Pillars (Q1-Q4)
			\item Three Lines of Defense — Roles and Responsibilities (Q5-Q8)
			\item Risk Appetite Statement (RAS) (Q9-Q12)
			\item Enterprise Risk Assessment — Inherent and Residual Risk (Q13-Q16)
			\item Risk-Based Approach (RBA) (Q17-Q20)
			\item Policies, Standards, and Procedures (Q21-Q24)
			\item Regulatory Guidance and Horizon Scanning (Q25-Q28)
			\item Governing Committees and Decision Making (Q29-Q32)
			\item KPIs, KRIs, and Management Reporting (Q33-Q36)
			\item FC Reporting Requirements — Jurisdictional Differences (Q37-Q40)
			\item Other Risk Management Teams (Q41-Q44)
			\item Controls at Onboarding and Throughout Lifecycle (Q45-Q48)
			\item Designing, Implementing, and Monitoring Controls (Q49-Q52)
			\item Risk Assessment — Inputs and Outputs (Q53-Q56)
			\item High-Risk Scenario Controls (Q57-Q60)
			\item Other Due Diligence — Employees, Vendors, Third Parties (Q61-Q64)
			\item Investigating Unusual Transactions (Q65-Q68)
			\item Documentation, Escalation, and SAR Writing (Q69-Q72)
			\item Governance Escalation and Offboarding (Q73-Q76)
			\item Responding to Court Orders and Law Enforcement (Q77-Q80)
			\item De-risking and Financial Inclusion (Q81-Q84)
			\item Customer Communication and Tipping-Off (Q85-Q88)
			\item Training and Awareness (Q89-Q92)
			\item AFC Team Structure and Roles (Q93-Q96)
			\item AFC and Front Office Interaction (Q97-Q100)
			\item Other Functions Supporting AFC (Q101-Q104)
		\end{enumerate}
	\end{frame}
	
	% ================================================================
	% SLIDE 3-6: EFFECTIVE AFC PROGRAM
	% ================================================================
	\begin{frame}
		\frametitle{\fontsize{11}{13}\selectfont Q1: Core Pillars of an AFC Program}
		\begin{block}{\fontsize{8}{10}\selectfont Topic: Effective AFC Program and Core Pillars}
			\scriptsize
			What are the core pillars of an effective AML/CFT compliance program under the US Bank Secrecy Act?
		\end{block}
		\vspace{0.05cm}
		\stepbox{\scriptsize
			\keypoint{Correct Answer:} Internal controls, designated compliance officer, employee training, independent testing, and CDD (including beneficial ownership) — the five pillars.
		}
		\vspace{0.05cm}
		\scriptsize
		\hardconcept{Core Concept:} The \textbf{BSA Five Pillars} (original four plus CDD Rule):
		\par \textbf{Pillar 1}: Internal controls — policies, procedures, and systems.
		\par \textbf{Pillar 2}: Designated compliance officer (BSA Officer/MLRO).
		\par \textbf{Pillar 3}: Ongoing employee training.
		\par \textbf{Pillar 4}: Independent testing/audit.
		\par \textbf{Pillar 5}: Customer Due Diligence (including beneficial ownership) — added 2016.
		\vspace{0.03cm}
		\wrongbox{\scriptsize \textbf{Common Trap:} "Only four pillars" — INCORRECT. The CDD Rule added a fifth pillar in 2016.}
		\vspace{0.03cm}
		\tipbox{\scriptsize \textbf{Key Insight:} All five pillars are \textbf{required} by the BSA. Regulators expect each to be implemented effectively.}
	\end{frame}
	
	\begin{frame}
		\frametitle{\fontsize{11}{13}\selectfont Q2: What Makes a Compliance Program Effective}
		\begin{block}{\fontsize{8}{10}\selectfont Topic: Effective Compliance Program}
			\scriptsize
			What makes an AML/CFT compliance program effective?
		\end{block}
		\vspace{0.05cm}
		\stepbox{\scriptsize
			\keypoint{Correct Answer:} A risk-based approach, clear accountability, adequate resources, and continuous monitoring and improvement.
		}
		\vspace{0.05cm}
		\scriptsize
		\hardconcept{Core Concept:} Key success factors:
		\par \textbf{Risk-Based}: Tailored to the institution's risk profile.
		\par \textbf{Accountability}: Clear roles and responsibilities.
		\par \textbf{Resources}: Adequate people, systems, and budget.
		\par \textbf{Continuous Improvement}: Regular review and updates.
		\par \textbf{Tone from the Top}: Senior management commitment.
		\vspace{0.03cm}
		\wrongbox{\scriptsize \textbf{Common Trap:} "Having written policies is enough" — INCORRECT. Effective programs must be \textbf{implemented and operational}.}
		\vspace{0.03cm}
		\tipbox{\scriptsize \textbf{Key Insight:} Regulators look for \textbf{effectiveness}, not just technical compliance.}
	\end{frame}
	
	\begin{frame}
		\frametitle{\fontsize{11}{13}\selectfont Q3: Compliance Program Resources}
		\begin{block}{\fontsize{8}{10}\selectfont Topic: Compliance Resources}
			\scriptsize
			What resources are needed for an effective AFC compliance program?
		\end{block}
		\vspace{0.05cm}
		\stepbox{\scriptsize
			\keypoint{Correct Answer:} Adequate staffing, appropriate technology systems, sufficient budget, and access to information.
		}
		\vspace{0.03cm}
		\tipbox{\scriptsize \textbf{Key Insight:} Under-resourced programs are a common regulatory finding.}
	\end{frame}
	
	\begin{frame}
		\frametitle{\fontsize{11}{13}\selectfont Q4: Tone from the Top}
		\begin{block}{\fontsize{8}{10}\selectfont Topic: Tone from the Top}
			\scriptsize
			Why is "tone from the top" important for an AFC program?
		\end{block}
		\vspace{0.05cm}
		\stepbox{\scriptsize
			\keypoint{Correct Answer:} Senior management commitment sets the cultural tone and demonstrates that compliance is a priority across the organization.
		}
		\vspace{0.03cm}
		\tipbox{\scriptsize \textbf{Key Insight:} Without tone from the top, \textbf{compliance culture} cannot take root.}
	\end{frame}
	
	% ================================================================
	% SLIDE 7-10: THREE LINES OF DEFENSE
	% ================================================================
	\begin{frame}
		\frametitle{\fontsize{11}{13}\selectfont Q5: Three Lines of Defense — Overview}
		\begin{block}{\fontsize{8}{10}\selectfont Topic: Three Lines of Defense}
			\scriptsize
			What are the three lines of defense in AML/CFT compliance?
		\end{block}
		\vspace{0.05cm}
		\stepbox{\scriptsize
			\keypoint{Correct Answer:} First line: Operational management (business units); Second line: Compliance and risk management (oversight); Third line: Internal audit (independent assurance).
		}
		\vspace{0.05cm}
		\scriptsize
		\hardconcept{Core Concept:}
		\par \textbf{First Line}: Business units identify, assess, and control risks.
		\par \textbf{Second Line}: AML compliance provides oversight and challenges the first line.
		\par \textbf{Third Line}: Internal audit provides independent assurance.
		\vspace{0.03cm}
		\wrongbox{\scriptsize \textbf{Common Trap:} "Compliance is the first line of defense" — INCORRECT. The first line is \textbf{operational management}, not compliance.}
		\vspace{0.03cm}
		\tipbox{\scriptsize \textbf{Key Insight:} The three lines must operate \textbf{independently but collaboratively}.}
	\end{frame}
	
	\begin{frame}
		\frametitle{\fontsize{11}{13}\selectfont Q6: First Line Responsibilities}
		\begin{block}{\fontsize{8}{10}\selectfont Topic: First Line Responsibilities}
			\scriptsize
			What are the responsibilities of the first line of defense in AML/CFT?
		\end{block}
		\vspace{0.05cm}
		\stepbox{\scriptsize
			\keypoint{Correct Answer:} Implementing AML controls in day-to-day operations, conducting CDD, identifying suspicious activity, and escalating concerns.
		}
		\vspace{0.03cm}
		\tipbox{\scriptsize \textbf{Key Insight:} The first line \textbf{owns} the risk and \textbf{implements} controls.}
	\end{frame}
	
	\begin{frame}
		\frametitle{\fontsize{11}{13}\selectfont Q7: Second Line Responsibilities}
		\begin{block}{\fontsize{8}{10}\selectfont Topic: Second Line Responsibilities}
			\scriptsize
			What are the responsibilities of the second line of defense?
		\end{block}
		\vspace{0.05cm}
		\stepbox{\scriptsize
			\keypoint{Correct Answer:} Overseeing and challenging the first line, developing AML policies, conducting risk assessments, and filing SARs.
		}
		\vspace{0.03cm}
		\tipbox{\scriptsize \textbf{Key Insight:} The second line \textbf{oversees} and \textbf{challenges} — it does not replace the first line.}
	\end{frame}
	
	\begin{frame}
		\frametitle{\fontsize{11}{13}\selectfont Q8: Third Line Responsibilities}
		\begin{block}{\fontsize{8}{10}\selectfont Topic: Third Line Responsibilities}
			\scriptsize
			What are the responsibilities of the third line of defense?
		\end{block}
		\vspace{0.05cm}
		\stepbox{\scriptsize
			\keypoint{Correct Answer:} Providing independent assurance by auditing the effectiveness of the first and second lines.
		}
		\vspace{0.03cm}
		\wrongbox{\scriptsize \textbf{Common Trap:} "Audit should design controls" — INCORRECT. Audit \textbf{cannot} design controls it later tests (independence violation).}
		\vspace{0.03cm}
		\tipbox{\scriptsize \textbf{Key Insight:} The third line must have \textbf{direct access} to the board.}
	\end{frame}
	
	% ================================================================
	% SLIDE 11-14: RISK APPETITE STATEMENT
	% ================================================================
	\begin{frame}
		\frametitle{\fontsize{11}{13}\selectfont Q9: Risk Appetite Statement — Purpose}
		\begin{block}{\fontsize{8}{10}\selectfont Topic: Risk Appetite Statement}
			\scriptsize
			What is the purpose of a Risk Appetite Statement (RAS) in AML/CFT?
		\end{block}
		\vspace{0.05cm}
		\stepbox{\scriptsize
			\keypoint{Correct Answer:} To define the level of risk the organization is willing to accept in pursuit of its business objectives, guiding decision-making and resource allocation.
		}
		\vspace{0.05cm}
		\scriptsize
		\hardconcept{Core Concept:} RAS components:
		\par \textbf{Strategic}: Aligns risk-taking with business strategy.
		\par \textbf{Qualitative}: Describes risk tolerance in words.
		\par \textbf{Quantitative}: Sets numerical limits where appropriate.
		\par \textbf{Communicated}: Clearly articulated to all stakeholders.
		\vspace{0.03cm}
		\tipbox{\scriptsize \textbf{Key Insight:} The RAS sets the \textbf{boundaries} for risk-taking in the organization.}
	\end{frame}
	
	\begin{frame}
		\frametitle{\fontsize{11}{13}\selectfont Q10: Defining a RAS}
		\begin{block}{\fontsize{8}{10}\selectfont Topic: Defining a RAS}
			\scriptsize
			How should a Risk Appetite Statement be defined?
		\end{block}
		\vspace{0.05cm}
		\stepbox{\scriptsize
			\keypoint{Correct Answer:} Through consultation with the board, senior management, and risk owners, based on the institution's risk assessment and strategic objectives.
		}
		\vspace{0.03cm}
		\tipbox{\scriptsize \textbf{Key Insight:} The RAS should be \textbf{approved by the board}.}
	\end{frame}
	
	\begin{frame}
		\frametitle{\fontsize{11}{13}\selectfont Q11: Communicating the RAS}
		\begin{block}{\fontsize{8}{10}\selectfont Topic: Communicating the RAS}
			\scriptsize
			How should a RAS be communicated in the organization?
		\end{block}
		\vspace{0.05cm}
		\stepbox{\scriptsize
			\keypoint{Correct Answer:} Through training, policy documents, and management communications to ensure all employees understand the organization's risk tolerance.
		}
		\vspace{0.03cm}
		\tipbox{\scriptsize \textbf{Key Insight:} Communication ensures the RAS is \textbf{operationalized}.}
	\end{frame}
	
	\begin{frame}
		\frametitle{\fontsize{11}{13}\selectfont Q12: Implementing Controls Based on RAS}
		\begin{block}{\fontsize{8}{10}\selectfont Topic: Controls and RAS}
			\scriptsize
			How does the RAS inform control design?
		\end{block}
		\vspace{0.05cm}
		\stepbox{\scriptsize
			\keypoint{Correct Answer:} Controls should be designed to ensure risk-taking remains within the defined appetite, with escalation for exceptions.
		}
		\vspace{0.03cm}
		\tipbox{\scriptsize \textbf{Key Insight:} The RAS is the \textbf{foundation} for control design.}
	\end{frame}
	
	% ================================================================
	% SLIDE 15-18: ENTERPRISE RISK ASSESSMENT
	% ================================================================
	\begin{frame}
		\frametitle{\fontsize{11}{13}\selectfont Q13: Enterprise Risk Assessment — Overview}
		\begin{block}{\fontsize{8}{10}\selectfont Topic: Enterprise Risk Assessment}
			\scriptsize
			What is an Enterprise-Wide Risk Assessment (EWRA)?
		\end{block}
		\vspace{0.05cm}
		\stepbox{\scriptsize
			\keypoint{Correct Answer:} A comprehensive assessment of the organization's ML/TF, sanctions, fraud, ABC, and tax evasion risks across all business lines and geographies.
		}
		\vspace{0.05cm}
		\scriptsize
		\hardconcept{Core Concept:} EWRA components:
		\par \textbf{Inherent Risk}: Risk before controls.
		\par \textbf{Controls}: Measures to mitigate risk.
		\par \textbf{Residual Risk}: Risk after controls.
		\par \textbf{Risk Assessment}: Evaluation of inherent and residual risk.
		\vspace{0.03cm}
		\tipbox{\scriptsize \textbf{Key Insight:} Under \textbf{FATF Recommendation 1}, institutions must conduct an EWRA.}
	\end{frame}
	
	\begin{frame}
		\frametitle{\fontsize{11}{13}\selectfont Q14: Identifying Inherent Risk}
		\begin{block}{\fontsize{8}{10}\selectfont Topic: Inherent Risk}
			\scriptsize
			How is inherent risk identified in an EWRA?
		\end{block}
		\vspace{0.05cm}
		\stepbox{\scriptsize
			\keypoint{Correct Answer:} By assessing the institution's products, services, customers, geographies, and channels without considering controls.
		}
		\vspace{0.03cm}
		\tipbox{\scriptsize \textbf{Key Insight:} Inherent risk is the \textbf{starting point} for the risk assessment.}
	\end{frame}
	
	\begin{frame}
		\frametitle{\fontsize{11}{13}\selectfont Q15: Controls and Residual Risk}
		\begin{block}{\fontsize{8}{10}\selectfont Topic: Residual Risk}
			\scriptsize
			How is residual risk assessed?
		\end{block}
		\vspace{0.05cm}
		\stepbox{\scriptsize
			\keypoint{Correct Answer:} By evaluating the effectiveness of controls in mitigating inherent risk and determining the remaining risk exposure.
		}
		\vspace{0.03cm}
		\tipbox{\scriptsize \textbf{Key Insight:} Residual risk must be \textbf{within the RAS}.}
	\end{frame}
	
	\begin{frame}
		\frametitle{\fontsize{11}{13}\selectfont Q16: Updating the EWRA}
		\begin{block}{\fontsize{8}{10}\selectfont Topic: EWRA Updates}
			\scriptsize
			When should the EWRA be updated?
		\end{block}
		\vspace{0.05cm}
		\stepbox{\scriptsize
			\keypoint{Correct Answer:} At least annually, and when there are significant changes to the business, regulatory landscape, or risk environment.
		}
		\vspace{0.03cm}
		\tipbox{\scriptsize \textbf{Key Insight:} The EWRA should be a \textbf{living document}.}
	\end{frame}
	
	% ================================================================
	% SLIDE 19-22: RISK-BASED APPROACH
	% ================================================================
	\begin{frame}
		\frametitle{\fontsize{11}{13}\selectfont Q17: Risk-Based Approach — Definition}
		\begin{block}{\fontsize{8}{10}\selectfont Topic: Risk-Based Approach}
			\scriptsize
			What is the risk-based approach in AML/CFT?
		\end{block}
		\vspace{0.05cm}
		\stepbox{\scriptsize
			\keypoint{Correct Answer:} Allocating resources and applying controls proportionately to the level of ML/TF risk identified.
		}
		\vspace{0.05cm}
		\scriptsize
		\hardconcept{Core Concept:} RBA principles:
		\par \textbf{Proportionality}: Controls match risk level.
		\par \textbf{Resource Allocation}: Focus on highest risks.
		\par \textbf{Flexibility}: Adapt to changing risks.
		\par \textbf{Justification}: Document risk-based decisions.
		\vspace{0.03cm}
		\wrongbox{\scriptsize \textbf{Common Trap:} "Risk-based approach means doing less compliance" — INCORRECT. It means \textbf{targeting} resources where they are most needed.}
		\vspace{0.03cm}
		\tipbox{\scriptsize \textbf{Key Insight:} \textbf{FATF Recommendation 1} establishes the RBA as the foundation of AML/CFT.}
	\end{frame}
	
	\begin{frame}
		\frametitle{\fontsize{11}{13}\selectfont Q18: Building Controls Commensurate with Risks}
		\begin{block}{\fontsize{8}{10}\selectfont Topic: Controls and RBA}
			\scriptsize
			How should controls be designed under the risk-based approach?
		\end{block}
		\vspace{0.05cm}
		\stepbox{\scriptsize
			\keypoint{Correct Answer:} Controls should be proportionate to the risk level — stronger controls for higher-risk activities, streamlined controls for lower-risk activities.
		}
		\vspace{0.03cm}
		\tipbox{\scriptsize \textbf{Key Insight:} \textbf{Proportionality} is the core of the RBA.}
	\end{frame}
	
	\begin{frame}
		\frametitle{\fontsize{11}{13}\selectfont Q19: Risk Assessment Frequency}
		\begin{block}{\fontsize{8}{10}\selectfont Topic: Risk Assessment Frequency}
			\scriptsize
			How often should risk assessments be conducted?
		\end{block}
		\vspace{0.05cm}
		\stepbox{\scriptsize
			\keypoint{Correct Answer:} Regularly (at least annually), and when there are significant changes in the business or regulatory environment.
		}
		\vspace{0.03cm}
		\tipbox{\scriptsize \textbf{Key Insight:} Continuous monitoring is also required.}
	\end{frame}
	
	\begin{frame}
		\frametitle{\fontsize{11}{13}\selectfont Q20: Documenting Risk-Based Decisions}
		\begin{block}{\fontsize{8}{10}\selectfont Topic: Documentation}
			\scriptsize
			Why is it important to document risk-based decisions?
		\end{block}
		\vspace{0.05cm}
		\stepbox{\scriptsize
			\keypoint{Correct Answer:} To demonstrate to regulators that the institution has considered risks and applied proportionate controls.
		}
		\vspace{0.03cm}
		\tipbox{\scriptsize \textbf{Key Insight:} Documentation provides \textbf{evidence} of compliance.}
	\end{frame}
	
	% ================================================================
	% SLIDE 23-26: POLICIES, STANDARDS, PROCEDURES
	% ================================================================
	\begin{frame}
		\frametitle{\fontsize{11}{13}\selectfont Q21: Policies vs Standards vs Procedures}
		\begin{block}{\fontsize{8}{10}\selectfont Topic: Policies, Standards, Procedures}
			\scriptsize
			What is the difference between policies, standards, and procedures in AML/CFT?
		\end{block}
		\vspace{0.05cm}
		\stepbox{\scriptsize
			\keypoint{Correct Answer:} Policies are high-level principles; standards are mandatory requirements; procedures are step-by-step instructions for implementation.
		}
		\vspace{0.05cm}
		\scriptsize
		\hardconcept{Core Concept:}
		\par \textbf{Policy}: "What" — high-level principles and commitments.
		\par \textbf{Standard}: "What must be done" — mandatory requirements.
		\par \textbf{Procedure}: "How" — step-by-step instructions.
		\vspace{0.03cm}
		\tipbox{\scriptsize \textbf{Key Insight:} All three are \textbf{necessary} for effective implementation.}
	\end{frame}
	
	\begin{frame}
		\frametitle{\fontsize{11}{13}\selectfont Q22: What Policies Should Contain}
		\begin{block}{\fontsize{8}{10}\selectfont Topic: Policy Content}
			\scriptsize
			What should AML/CFT policies contain?
		\end{block}
		\vspace{0.05cm}
		\stepbox{\scriptsize
			\keypoint{Correct Answer:} Purpose, scope, roles and responsibilities, risk assessment approach, control requirements, and reporting/escalation procedures.
		}
		\vspace{0.03cm}
		\tipbox{\scriptsize \textbf{Key Insight:} Policies should be \textbf{approved by senior management}.}
	\end{frame}
	
	\begin{frame}
		\frametitle{\fontsize{11}{13}\selectfont Q23: Standards in AML/CFT}
		\begin{block}{\fontsize{8}{10}\selectfont Topic: Standards}
			\scriptsize
			What are standards in the AML/CFT context?
		\end{block}
		\vspace{0.05cm}
		\stepbox{\scriptsize
			\keypoint{Correct Answer:} Mandatory requirements that must be met to comply with policies (e.g., CDD standards, transaction monitoring thresholds).
		}
		\vspace{0.03cm}
		\tipbox{\scriptsize \textbf{Key Insight:} Standards are \textbf{measurable} and \textbf{enforceable}.}
	\end{frame}
	
	\begin{frame}
		\frametitle{\fontsize{11}{13}\selectfont Q24: Procedures in AML/CFT}
		\begin{block}{\fontsize{8}{10}\selectfont Topic: Procedures}
			\scriptsize
			What are procedures in the AML/CFT context?
		\end{block}
		\vspace{0.05cm}
		\stepbox{\scriptsize
			\keypoint{Correct Answer:} Detailed, step-by-step instructions for how to implement policies and standards in day-to-day operations.
		}
		\vspace{0.03cm}
		\tipbox{\scriptsize \textbf{Key Insight:} Procedures should be \textbf{clear} and \textbf{accessible}.}
	\end{frame}
	
	% ================================================================
	% SLIDE 27-30: REGULATORY GUIDANCE AND HORIZON SCANNING
	% ================================================================
	\begin{frame}
		\frametitle{\fontsize{11}{13}\selectfont Q25: Reviewing Regulatory Guidance}
		\begin{block}{\fontsize{8}{10}\selectfont Topic: Regulatory Guidance}
			\scriptsize
			Why should institutions review regulatory guidance when drafting policies?
		\end{block}
		\vspace{0.05cm}
		\stepbox{\scriptsize
			\keypoint{Correct Answer:} To ensure policies align with regulatory expectations and incorporate the latest requirements.
		}
		\vspace{0.03cm}
		\tipbox{\scriptsize \textbf{Key Insight:} Regulatory guidance provides \textbf{clarity} on expectations.}
	\end{frame}
	
	\begin{frame}
		\frametitle{\fontsize{11}{13}\selectfont Q26: Horizon Scanning}
		\begin{block}{\fontsize{8}{10}\selectfont Topic: Horizon Scanning}
			\scriptsize
			What is horizon scanning in AML/CFT?
		\end{block}
		\vspace{0.05cm}
		\stepbox{\scriptsize
			\keypoint{Correct Answer:} The process of identifying and anticipating emerging regulatory, technological, and criminal trends that may impact the institution.
		}
		\vspace{0.03cm}
		\tipbox{\scriptsize \textbf{Key Insight:} Horizon scanning helps institutions \textbf{prepare} for future changes.}
	\end{frame}
	
	\begin{frame}
		\frametitle{\fontsize{11}{13}\selectfont Q27: Incorporating Industry Best Practices}
		\begin{block}{\fontsize{8}{10}\selectfont Topic: Best Practices}
			\scriptsize
			Why should institutions incorporate industry best practices?
		\end{block}
		\vspace{0.05cm}
		\stepbox{\scriptsize
			\keypoint{Correct Answer:} To enhance the effectiveness of the compliance program and meet regulatory expectations.
		}
		\vspace{0.03cm}
		\tipbox{\scriptsize \textbf{Key Insight:} Regulators often reference industry best practices as benchmarks.}
	\end{frame}
	
	\begin{frame}
		\frametitle{\fontsize{11}{13}\selectfont Q28: Sources of Guidance}
		\begin{block}{\fontsize{8}{10}\selectfont Topic: Guidance Sources}
			\scriptsize
			What are sources of regulatory guidance for AML/CFT?
		\end{block}
		\vspace{0.05cm}
		\stepbox{\scriptsize
			\keypoint{Correct Answer:} FATF Recommendations, national regulators (FinCEN, FCA), FSRBs, and industry bodies (Wolfsberg Group, Basel Committee).
		}
		\vspace{0.03cm}
		\tipbox{\scriptsize \textbf{Key Insight:} Use \textbf{multiple sources} for comprehensive guidance.}
	\end{frame}
	
	% ================================================================
	% SLIDE 31-34: GOVERNING COMMITTEES
	% ================================================================
	\begin{frame}
		\frametitle{\fontsize{11}{13}\selectfont Q29: Governing Committees — Purpose}
		\begin{block}{\fontsize{8}{10}\selectfont Topic: Governing Committees}
			\scriptsize
			What is the purpose of governing committees in AML/CFT?
		\end{block}
		\vspace{0.05cm}
		\stepbox{\scriptsize
			\keypoint{Correct Answer:} To provide oversight, make decisions on risk acceptance, and ensure the compliance program is effective.
		}
		\vspace{0.03cm}
		\tipbox{\scriptsize \textbf{Key Insight:} Committees provide \textbf{accountability} and \textbf{decision-making}.}
	\end{frame}
	
	\begin{frame}
		\frametitle{\fontsize{11}{13}\selectfont Q30: Committee Structure}
		\begin{block}{\fontsize{8}{10}\selectfont Topic: Committee Structure}
			\scriptsize
			What should the structure of governing committees include?
		\end{block}
		\vspace{0.05cm}
		\stepbox{\scriptsize
			\keypoint{Correct Answer:} Clear terms of reference, defined membership, regular meetings, and documented decisions.
		}
		\vspace{0.03cm}
		\tipbox{\scriptsize \textbf{Key Insight:} Committees should have \textbf{terms of reference} that define their authority.}
	\end{frame}
	
	\begin{frame}
		\frametitle{\fontsize{11}{13}\selectfont Q31: Committee Decision Making}
		\begin{block}{\fontsize{8}{10}\selectfont Topic: Committee Decisions}
			\scriptsize
			What decisions should governing committees make?
		\end{block}
		\vspace{0.05cm}
		\stepbox{\scriptsize
			\keypoint{Correct Answer:} Risk acceptance, policy approvals, escalation decisions, and oversight of the compliance program.
		}
		\vspace{0.03cm}
		\tipbox{\scriptsize \textbf{Key Insight:} Major decisions should be \textbf{escalated} to the committee.}
	\end{frame}
	
	\begin{frame}
		\frametitle{\fontsize{11}{13}\selectfont Q32: Committee Documentation}
		\begin{block}{\fontsize{8}{10}\selectfont Topic: Committee Documentation}
			\scriptsize
			Why is committee documentation important?
		\end{block}
		\vspace{0.05cm}
		\stepbox{\scriptsize
			\keypoint{Correct Answer:} To provide an audit trail of decisions and demonstrate oversight.
		}
		\vspace{0.03cm}
		\tipbox{\scriptsize \textbf{Key Insight:} \textbf{Minutes} should be recorded for all committee meetings.}
	\end{frame}
	
	% ================================================================
	% SLIDE 35-38: KPIs, KRIs, MANAGEMENT REPORTING
	% ================================================================
	\begin{frame}
		\frametitle{\fontsize{11}{13}\selectfont Q33: KPIs vs KRIs}
		\begin{block}{\fontsize{8}{10}\selectfont Topic: KPIs and KRIs}
			\scriptsize
			What is the difference between KPIs and KRIs in AML/CFT?
		\end{block}
		\vspace{0.05cm}
		\stepbox{\scriptsize
			\keypoint{Correct Answer:} KPIs measure performance (e.g., alerts reviewed), while KRIs measure risk exposure (e.g., number of high-risk customers).
		}
		\vspace{0.05cm}
		\scriptsize
		\hardconcept{Core Concept:}
		\par \textbf{KPI (Key Performance Indicator)}: Measures \textbf{effectiveness} and \textbf{efficiency}.
		\par \textbf{KRI (Key Risk Indicator)}: Measures \textbf{risk exposure} and \textbf{trends}.
		\vspace{0.03cm}
		\tipbox{\scriptsize \textbf{Key Insight:} Both are needed for comprehensive management reporting.}
	\end{frame}
	
	\begin{frame}
		\frametitle{\fontsize{11}{13}\selectfont Q34: Management Reporting Information}
		\begin{block}{\fontsize{8}{10}\selectfont Topic: Management Reporting}
			\scriptsize
			What information should be included in management reporting?
		\end{block}
		\vspace{0.05cm}
		\stepbox{\scriptsize
			\keypoint{Correct Answer:} KPI/KRI trends, emerging risks, regulatory changes, audit findings, and material incidents.
		}
		\vspace{0.03cm}
		\tipbox{\scriptsize \textbf{Key Insight:} Reports should be \textbf{regular}, \textbf{clear}, and \textbf{actionable}.}
	\end{frame}
	
	\begin{frame}
		\frametitle{\fontsize{11}{13}\selectfont Q35: Emerging Risk Reporting}
		\begin{block}{\fontsize{8}{10}\selectfont Topic: Emerging Risks}
			\scriptsize
			How should emerging financial crime risks be reported?
		\end{block}
		\vspace{0.05cm}
		\stepbox{\scriptsize
			\keypoint{Correct Answer:} Through regular horizon scanning reports and ad-hoc updates to management and the board.
		}
		\vspace{0.03cm}
		\tipbox{\scriptsize \textbf{Key Insight:} Emerging risks should be \textbf{escalated} promptly.}
	\end{frame}
	
	\begin{frame}
		\frametitle{\fontsize{11}{13}\selectfont Q36: Board Oversight Information}
		\begin{block}{\fontsize{8}{10}\selectfont Topic: Board Reporting}
			\scriptsize
			What information should be provided to the board?
		\end{block}
		\vspace{0.05cm}
		\stepbox{\scriptsize
			\keypoint{Correct Answer:} Overall program effectiveness, significant risks, regulatory issues, and compliance with the RAS.
		}
		\vspace{0.03cm}
		\tipbox{\scriptsize \textbf{Key Insight:} The board needs \textbf{strategic} reporting, not just operational details.}
	\end{frame}
	
	% ================================================================
	% SLIDE 39-42: FC REPORTING REQUIREMENTS
	% ================================================================
	\begin{frame}
		\frametitle{\fontsize{11}{13}\selectfont Q37: STR Filing Requirements}
		\begin{block}{\fontsize{8}{10}\selectfont Topic: FC Reporting}
			\scriptsize
			What is a Suspicious Activity Report (SAR) filing requirement?
		\end{block}
		\vspace{0.05cm}
		\stepbox{\scriptsize
			\keypoint{Correct Answer:} Financial institutions must file SARs when they suspect money laundering, terrorist financing, or other financial crimes.
		}
		\vspace{0.03cm}
		\tipbox{\scriptsize \textbf{Key Insight:} Under \textbf{FATF Recommendation 20}, suspicious transactions must be reported.}
	\end{frame}
	
	\begin{frame}
		\frametitle{\fontsize{11}{13}\selectfont Q38: Jurisdictional Reporting Differences}
		\begin{block}{\fontsize{8}{10}\selectfont Topic: Jurisdictional Differences}
			\scriptsize
			How do reporting requirements differ across jurisdictions?
		\end{block}
		\vspace{0.05cm}
		\stepbox{\scriptsize
			\keypoint{Correct Answer:} Different thresholds, reporting formats, filing deadlines, and competent authorities (e.g., FinCEN in US, NCA in UK).
		}
		\vspace{0.03cm}
		\tipbox{\scriptsize \textbf{Key Insight:} Global institutions must comply with \textbf{multiple} reporting regimes.}
	\end{frame}
	
	\begin{frame}
		\frametitle{\fontsize{11}{13}\selectfont Q39: Currency Transaction Reporting}
		\begin{block}{\fontsize{8}{10}\selectfont Topic: CTR}
			\scriptsize
			What is a Currency Transaction Report (CTR)?
		\end{block}
		\vspace{0.05cm}
		\stepbox{\scriptsize
			\keypoint{Correct Answer:} A report filed for cash transactions exceeding a threshold (e.g., \$10,000 in the US).
		}
		\vspace{0.03cm}
		\tipbox{\scriptsize \textbf{Key Insight:} CTRs are separate from SARs and have different thresholds.}
	\end{frame}
	
	\begin{frame}
		\frametitle{\fontsize{11}{13}\selectfont Q40: Reporting Timelines}
		\begin{block}{\fontsize{8}{10}\selectfont Topic: Reporting Timelines}
			\scriptsize
			What are typical reporting timelines for SARs?
		\end{block}
		\vspace{0.05cm}
		\stepbox{\scriptsize
			\keypoint{Correct Answer:} Generally 15-30 days from detection, or immediately for urgent cases.
		}
		\vspace{0.03cm}
		\tipbox{\scriptsize \textbf{Key Insight:} \textbf{Timely} reporting is a regulatory expectation.}
	\end{frame}
	
	% ================================================================
	% SLIDE 43-46: OTHER RISK MANAGEMENT TEAMS
	% ================================================================
	\begin{frame}
		\frametitle{\fontsize{11}{13}\selectfont Q41: Operational Risk and AFC}
		\begin{block}{\fontsize{8}{10}\selectfont Topic: Operational Risk}
			\scriptsize
			How does operational risk relate to AFC compliance?
		\end{block}
		\vspace{0.05cm}
		\stepbox{\scriptsize
			\keypoint{Correct Answer:} AFC failures are operational risk events that can result in financial loss, regulatory penalties, and reputational damage.
		}
		\vspace{0.03cm}
		\tipbox{\scriptsize \textbf{Key Insight:} AFC is a subset of \textbf{operational risk}.}
	\end{frame}
	
	\begin{frame}
		\frametitle{\fontsize{11}{13}\selectfont Q42: Financial Risk and AFC}
		\begin{block}{\fontsize{8}{10}\selectfont Topic: Financial Risk}
			\scriptsize
			How does financial risk relate to AFC compliance?
		\end{block}
		\vspace{0.05cm}
		\stepbox{\scriptsize
			\keypoint{Correct Answer:} AFC failures can cause direct financial losses through fines, penalties, and asset seizures.
		}
		\vspace{0.03cm}
		\tipbox{\scriptsize \textbf{Key Insight:} Financial risk includes \textbf{regulatory} and \textbf{reputational} components.}
	\end{frame}
	
	\begin{frame}
		\frametitle{\fontsize{11}{13}\selectfont Q43: Non-Financial Risk and AFC}
		\begin{block}{\fontsize{8}{10}\selectfont Topic: Non-Financial Risk}
			\scriptsize
			What non-financial risks are associated with AFC?
		\end{block}
		\vspace{0.05cm}
		\stepbox{\scriptsize
			\keypoint{Correct Answer:} Reputational damage, regulatory sanctions, and loss of license or business opportunities.
		}
		\vspace{0.03cm}
		\tipbox{\scriptsize \textbf{Key Insight:} Non-financial risks can be \textbf{more severe} than financial penalties.}
	\end{frame}
	
	\begin{frame}
		\frametitle{\fontsize{11}{13}\selectfont Q44: Reputational Risk and AFC}
		\begin{block}{\fontsize{8}{10}\selectfont Topic: Reputational Risk}
			\scriptsize
			How does AFC impact reputational risk?
		\end{block}
		\vspace{0.05cm}
		\stepbox{\scriptsize
			\keypoint{Correct Answer:} AFC failures can cause significant reputational damage, leading to loss of customers, partners, and market confidence.
		}
		\vspace{0.03cm}
		\tipbox{\scriptsize \textbf{Key Insight:} Reputational risk is a \textbf{top concern} for boards.}
	\end{frame}
	
	% ================================================================
	% SLIDE 47-50: CONTROLS AT ONBOARDING AND LIFECYCLE
	% ================================================================
	\begin{frame}
		\frametitle{\fontsize{11}{13}\selectfont Q45: Onboarding Controls}
		\begin{block}{\fontsize{8}{10}\selectfont Topic: Onboarding Controls}
			\scriptsize
			What controls should be applied at customer onboarding?
		\end{block}
		\vspace{0.05cm}
		\stepbox{\scriptsize
			\keypoint{Correct Answer:} CDD, identity verification, sanctions screening, PEP screening, and risk assessment.
		}
		\vspace{0.03cm}
		\tipbox{\scriptsize \textbf{Key Insight:} Onboarding is the \textbf{critical first step} in AFC compliance.}
	\end{frame}
	
	\begin{frame}
		\frametitle{\fontsize{11}{13}\selectfont Q46: Ongoing Monitoring Controls}
		\begin{block}{\fontsize{8}{10}\selectfont Topic: Ongoing Monitoring}
			\scriptsize
			What controls should be applied throughout the customer lifecycle?
		\end{block}
		\vspace{0.05cm}
		\stepbox{\scriptsize
			\keypoint{Correct Answer:} Transaction monitoring, periodic reviews, trigger-based reviews, and ongoing sanctions/PEP screening.
		}
		\vspace{0.03cm}
		\tipbox{\scriptsize \textbf{Key Insight:} Monitoring is \textbf{continuous}, not just at onboarding.}
	\end{frame}
	
	\begin{frame}
		\frametitle{\fontsize{11}{13}\selectfont Q47: Trigger-Based Reviews}
		\begin{block}{\fontsize{8}{10}\selectfont Topic: Trigger-Based Reviews}
			\scriptsize
			What triggers a customer review during the lifecycle?
		\end{block}
		\vspace{0.05cm}
		\stepbox{\scriptsize
			\keypoint{Correct Answer:} Significant changes in behavior, new risk factors, regulatory updates, or system alerts.
		}
		\vspace{0.03cm}
		\tipbox{\scriptsize \textbf{Key Insight:} Trigger-based reviews are \textbf{event-driven}.}
	\end{frame}
	
	\begin{frame}
		\frametitle{\fontsize{11}{13}\selectfont Q48: Exit Controls}
		\begin{block}{\fontsize{8}{10}\selectfont Topic: Exit Controls}
			\scriptsize
			What controls should be applied when exiting a customer relationship?
		\end{block}
		\vspace{0.05cm}
		\stepbox{\scriptsize
			\keypoint{Correct Answer:} Consideration at the relevant committee, liaison with law enforcement and FIUs regarding keeping accounts open as needed, and appropriate documentation.
		}
		\vspace{0.03cm}
		\tipbox{\scriptsize \textbf{Key Insight:} Exits should be \textbf{managed} and \textbf{escalated}.}
	\end{frame}
	
	% ================================================================
	% SLIDE 51-54: DESIGNING, IMPLEMENTING, MONITORING CONTROLS
	% ================================================================
	\begin{frame}
		\frametitle{\fontsize{11}{13}\selectfont Q49: KYC and CDD}
		\begin{block}{\fontsize{8}{10}\selectfont Topic: KYC and CDD}
			\scriptsize
			What is the difference between KYC and CDD?
		\end{block}
		\vspace{0.05cm}
		\stepbox{\scriptsize
			\keypoint{Correct Answer:} KYC is the process of identifying the customer; CDD includes KYC plus ongoing monitoring and risk assessment.
		}
		\vspace{0.03cm}
		\tipbox{\scriptsize \textbf{Key Insight:} CDD is a \textbf{broader} concept than KYC.}
	\end{frame}
	
	\begin{frame}
		\frametitle{\fontsize{11}{13}\selectfont Q50: Enhanced Due Diligence}
		\begin{block}{\fontsize{8}{10}\selectfont Topic: EDD}
			\scriptsize
			When is Enhanced Due Diligence (EDD) required?
		\end{block}
		\vspace{0.05cm}
		\stepbox{\scriptsize
			\keypoint{Correct Answer:} For high-risk customers, including PEPs, customers from high-risk jurisdictions, and those with complex ownership structures.
		}
		\vspace{0.03cm}
		\tipbox{\scriptsize \textbf{Key Insight:} EDD is \textbf{mandatory} for PEPs under FATF Recommendation 12.}
	\end{frame}
	
	\begin{frame}
		\frametitle{\fontsize{11}{13}\selectfont Q51: Employee Due Diligence}
		\begin{block}{\fontsize{8}{10}\selectfont Topic: Employee Due Diligence}
			\scriptsize
			What is employee due diligence in the AFC context?
		\end{block}
		\vspace{0.05cm}
		\stepbox{\scriptsize
			\keypoint{Correct Answer:} Background checks on employees, particularly those in sensitive roles, to identify potential insider threats.
		}
		\vspace{0.03cm}
		\tipbox{\scriptsize \textbf{Key Insight:} Insider threats are a significant \textbf{vulnerability}.}
	\end{frame}
	
	\begin{frame}
		\frametitle{\fontsize{11}{13}\selectfont Q52: Customer Risk Assessment}
		\begin{block}{\fontsize{8}{10}\selectfont Topic: Customer Risk Assessment}
			\scriptsize
			What factors should be considered in customer risk assessment?
		\end{block}
		\vspace{0.05cm}
		\stepbox{\scriptsize
			\keypoint{Correct Answer:} Customer type, geography, product usage, transaction patterns, and ownership structure.
		}
		\vspace{0.03cm}
		\tipbox{\scriptsize \textbf{Key Insight:} Risk assessment is \textbf{dynamic} and should be updated.}
	\end{frame}
	
	% ================================================================
	% SLIDE 55-58: RISK ASSESSMENT — INPUTS AND OUTPUTS
	% ================================================================
	\begin{frame}
		\frametitle{\fontsize{11}{13}\selectfont Q53: Types of Risk Assessment}
		\begin{block}{\fontsize{8}{10}\selectfont Topic: Types of Risk Assessment}
			\scriptsize
			What types of risk assessments are used in AFC?
		\end{block}
		\vspace{0.05cm}
		\stepbox{\scriptsize
			\keypoint{Correct Answer:} Enterprise-wide, business unit, product, geographic, customer, and transaction risk assessments.
		}
		\vspace{0.03cm}
		\tipbox{\scriptsize \textbf{Key Insight:} Different assessments serve \textbf{different purposes}.}
	\end{frame}
	
	\begin{frame}
		\frametitle{\fontsize{11}{13}\selectfont Q54: Inputs to Risk Assessment}
		\begin{block}{\fontsize{8}{10}\selectfont Topic: Risk Assessment Inputs}
			\scriptsize
			What are inputs to a risk assessment?
		\end{block}
		\vspace{0.05cm}
		\stepbox{\scriptsize
			\keypoint{Correct Answer:} National risk assessments, sectoral risk assessments, regulatory guidance, internal data, and threat intelligence.
		}
		\vspace{0.03cm}
		\tipbox{\scriptsize \textbf{Key Insight:} Inputs should be \textbf{diverse} and \textbf{current}.}
	\end{frame}
	
	\begin{frame}
		\frametitle{\fontsize{11}{13}\selectfont Q55: Outputs of Risk Assessment}
		\begin{block}{\fontsize{8}{10}\selectfont Topic: Risk Assessment Outputs}
			\scriptsize
			What are outputs of a risk assessment?
		\end{block}
		\vspace{0.05cm}
		\stepbox{\scriptsize
			\keypoint{Correct Answer:} Risk ratings, control recommendations, resource allocation decisions, and prioritization of activities.
		}
		\vspace{0.03cm}
		\tipbox{\scriptsize \textbf{Key Insight:} Outputs should be \textbf{actionable}.}
	\end{frame}
	
	\begin{frame}
		\frametitle{\fontsize{11}{13}\selectfont Q56: Managing AFC Risks}
		\begin{block}{\fontsize{8}{10}\selectfont Topic: Managing Risks}
			\scriptsize
			How are AFC risks managed after assessment?
		\end{block}
		\vspace{0.05cm}
		\stepbox{\scriptsize
			\keypoint{Correct Answer:} Through controls, monitoring, and regular reassessment to ensure risks remain within appetite.
		}
		\vspace{0.03cm}
		\tipbox{\scriptsize \textbf{Key Insight:} Risk management is a \textbf{continuous process}.}
	\end{frame}
	
	% ================================================================
	% SLIDE 59-62: HIGH-RISK SCENARIO CONTROLS
	% ================================================================
	\begin{frame}
		\frametitle{\fontsize{11}{13}\selectfont Q57: Conflict Zone Transactions}
		\begin{block}{\fontsize{8}{10}\selectfont Topic: Conflict Zone Controls}
			\scriptsize
			What controls should be applied to transactions involving conflict zones?
		\end{block}
		\vspace{0.05cm}
		\stepbox{\scriptsize
			\keypoint{Correct Answer:} Enhanced due diligence, sanctions screening, transaction monitoring, and escalation for approval.
		}
		\vspace{0.03cm}
		\tipbox{\scriptsize \textbf{Key Insight:} Conflict zones are \textbf{high-risk} jurisdictions.}
	\end{frame}
	
	\begin{frame}
		\frametitle{\fontsize{11}{13}\selectfont Q58: Dual-Use Goods}
		\begin{block}{\fontsize{8}{10}\selectfont Topic: Dual-Use Goods}
			\scriptsize
			What controls should be applied to dual-use goods transactions?
		\end{block}
		\vspace{0.05cm}
		\stepbox{\scriptsize
			\keypoint{Correct Answer:} Export control screening, end-use verification, and sanctions checks on parties.
		}
		\vspace{0.03cm}
		\tipbox{\scriptsize \textbf{Key Insight:} Dual-use goods can be diverted for \textbf{military} purposes.}
	\end{frame}
	
	\begin{frame}
		\frametitle{\fontsize{11}{13}\selectfont Q59: Complex Ownership Structures}
		\begin{block}{\fontsize{8}{10}\selectfont Topic: Complex Ownership}
			\scriptsize
			What controls should be applied to complex ownership structures?
		\end{block}
		\vspace{0.05cm}
		\stepbox{\scriptsize
			\keypoint{Correct Answer:} Identification of beneficial owners, source of funds verification, and enhanced monitoring.
		}
		\vspace{0.03cm}
		\tipbox{\scriptsize \textbf{Key Insight:} Complex structures are a \textbf{red flag} for ML.}
	\end{frame}
	
	\begin{frame}
		\frametitle{\fontsize{11}{13}\selectfont Q60: PEP Controls}
		\begin{block}{\fontsize{8}{10}\selectfont Topic: PEP Controls}
			\scriptsize
			What controls should be applied to PEPs?
		\end{block}
		\vspace{0.05cm}
		\stepbox{\scriptsize
			\keypoint{Correct Answer:} EDD, senior management approval, source of wealth verification, and ongoing monitoring.
		}
		\vspace{0.03cm}
		\tipbox{\scriptsize \textbf{Key Insight:} PEP controls are \textbf{mandatory} under FATF.}
	\end{frame}
	
	% ================================================================
	% SLIDE 63-66: OTHER DUE DILIGENCE
	% ================================================================
	\begin{frame}
		\frametitle{\fontsize{11}{13}\selectfont Q61: Employee Due Diligence}
		\begin{block}{\fontsize{8}{10}\selectfont Topic: Employee Due Diligence}
			\scriptsize
			Why is employee due diligence important in AFC?
		\end{block}
		\vspace{0.05cm}
		\stepbox{\scriptsize
			\keypoint{Correct Answer:} To mitigate insider threats and ensure employees in sensitive roles are trustworthy.
		}
		\vspace{0.03cm}
		\tipbox{\scriptsize \textbf{Key Insight:} Insider threats are a \textbf{significant vulnerability}.}
	\end{frame}
	
	\begin{frame}
		\frametitle{\fontsize{11}{13}\selectfont Q62: Vendor Due Diligence}
		\begin{block}{\fontsize{8}{10}\selectfont Topic: Vendor Due Diligence}
			\scriptsize
			What due diligence should be conducted on vendors?
		\end{block}
		\vspace{0.05cm}
		\stepbox{\scriptsize
			\keypoint{Correct Answer:} Background checks, sanctions screening, and assessment of their AML controls.
		}
		\vspace{0.03cm}
		\tipbox{\scriptsize \textbf{Key Insight:} Vendors can be a \textbf{weak link} in compliance.}
	\end{frame}
	
	\begin{frame}
		\frametitle{\fontsize{11}{13}\selectfont Q63: Third-Party Due Diligence}
		\begin{block}{\fontsize{8}{10}\selectfont Topic: Third-Party Due Diligence}
			\scriptsize
			What should third-party due diligence include?
		\end{block}
		\vspace{0.05cm}
		\stepbox{\scriptsize
			\keypoint{Correct Answer:} Assessment of the third party's AML program, sanctions screening, and contractual compliance obligations.
		}
		\vspace{0.03cm}
		\tipbox{\scriptsize \textbf{Key Insight:} Third parties can be used to \textbf{bypass} controls.}
	\end{frame}
	
	\begin{frame}
		\frametitle{\fontsize{11}{13}\selectfont Q64: Insider Threat Mitigation}
		\begin{block}{\fontsize{8}{10}\selectfont Topic: Insider Threats}
			\scriptsize
			How can insider threats be mitigated?
		\end{block}
		\vspace{0.05cm}
		\stepbox{\scriptsize
			\keypoint{Correct Answer:} Background checks, access controls, monitoring, and employee training.
		}
		\vspace{0.03cm}
		\tipbox{\scriptsize \textbf{Key Insight:} \textbf{Access controls} and \textbf{monitoring} are key controls.}
	\end{frame}
	
	% ================================================================
	% SLIDE 67-70: INVESTIGATING UNUSUAL TRANSACTIONS
	% ================================================================
	\begin{frame}
		\frametitle{\fontsize{11}{13}\selectfont Q65: Identifying Suspicious Activity}
		\begin{block}{\fontsize{8}{10}\selectfont Topic: Identifying Suspicious Activity}
			\scriptsize
			How can suspicious activity be identified?
		\end{block}
		\vspace{0.05cm}
		\stepbox{\scriptsize
			\keypoint{Correct Answer:} Through transaction monitoring alerts, staff reports, customer complaints, and external sources.
		}
		\vspace{0.03cm}
		\tipbox{\scriptsize \textbf{Key Insight:} \textbf{Multiple channels} are needed for detection.}
	\end{frame}
	
	\begin{frame}
		\frametitle{\fontsize{11}{13}\selectfont Q66: Staff Reports of Suspicion}
		\begin{block}{\fontsize{8}{10}\selectfont Topic: Staff Reports}
			\scriptsize
			Why are staff reports important in detection?
		\end{block}
		\vspace{0.05cm}
		\stepbox{\scriptsize
			\keypoint{Correct Answer:} Frontline staff may observe suspicious behavior that is not captured by automated systems.
		}
		\vspace{0.03cm}
		\tipbox{\scriptsize \textbf{Key Insight:} Human observation is a \textbf{critical control}.}
	\end{frame}
	
	\begin{frame}
		\frametitle{\fontsize{11}{13}\selectfont Q67: External Sources of Suspicion}
		\begin{block}{\fontsize{8}{10}\selectfont Topic: External Sources}
			\scriptsize
			What external sources can indicate suspicious activity?
		\end{block}
		\vspace{0.05cm}
		\stepbox{\scriptsize
			\keypoint{Correct Answer:} Law enforcement notifications, regulatory alerts, and adverse media.
		}
		\vspace{0.03cm}
		\tipbox{\scriptsize \textbf{Key Insight:} External sources provide \textbf{context} for alerts.}
	\end{frame}
	
	\begin{frame}
		\frametitle{\fontsize{11}{13}\selectfont Q68: Information Gathering}
		\begin{block}{\fontsize{8}{10}\selectfont Topic: Information Gathering}
			\scriptsize
			What information should be gathered during an investigation?
		\end{block}
		\vspace{0.05cm}
		\stepbox{\scriptsize
			\keypoint{Correct Answer:} Transaction details, customer profile, counterparty information, and source of funds.
		}
		\vspace{0.03cm}
		\tipbox{\scriptsize \textbf{Key Insight:} Investigations should be \textbf{thorough} and \textbf{documentation-based}.}
	\end{frame}
	
	% ================================================================
	% SLIDE 71-74: DOCUMENTATION AND ESCALATION
	% ================================================================
	\begin{frame}
		\frametitle{\fontsize{11}{13}\selectfont Q69: Investigation Documentation}
		\begin{block}{\fontsize{8}{10}\selectfont Topic: Investigation Documentation}
			\scriptsize
			Why is investigation documentation important?
		\end{block}
		\vspace{0.05cm}
		\stepbox{\scriptsize
			\keypoint{Correct Answer:} To provide an audit trail, support decision-making, and demonstrate compliance.
		}
		\vspace{0.03cm}
		\tipbox{\scriptsize \textbf{Key Insight:} Documentation should be \textbf{contemporaneous}.}
	\end{frame}
	
	\begin{frame}
		\frametitle{\fontsize{11}{13}\selectfont Q70: Writing SARs}
		\begin{block}{\fontsize{8}{10}\selectfont Topic: SAR Writing}
			\scriptsize
			What should a SAR include?
		\end{block}
		\vspace{0.05cm}
		\stepbox{\scriptsize
			\keypoint{Correct Answer:} Details of the suspicious activity, parties involved, amounts, and the basis for suspicion.
		}
		\vspace{0.03cm}
		\tipbox{\scriptsize \textbf{Key Insight:} SARs should be \textbf{clear}, \textbf{concise}, and \textbf{complete}.}
	\end{frame}
	
	\begin{frame}
		\frametitle{\fontsize{11}{13}\selectfont Q71: Internal Investigation Timelines}
		\begin{block}{\fontsize{8}{10}\selectfont Topic: Investigation Timelines}
			\scriptsize
			What are typical internal investigation timelines?
		\end{block}
		\vspace{0.05cm}
		\stepbox{\scriptsize
			\keypoint{Correct Answer:} Investigations should be completed within regulatory deadlines (e.g., 30 days for SAR filing in many jurisdictions).
		}
		\vspace{0.03cm}
		\tipbox{\scriptsize \textbf{Key Insight:} \textbf{Timeliness} is a regulatory expectation.}
	\end{frame}
	
	\begin{frame}
		\frametitle{\fontsize{11}{13}\selectfont Q72: Escalation Triggers}
		\begin{block}{\fontsize{8}{10}\selectfont Topic: Escalation}
			\scriptsize
			What triggers escalation of an investigation?
		\end{block}
		\vspace{0.05cm}
		\stepbox{\scriptsize
			\keypoint{Correct Answer:} Confirmed suspicion, significant risk, reputational concerns, or material thresholds.
		}
		\vspace{0.03cm}
		\tipbox{\scriptsize \textbf{Key Insight:} Escalation should be \textbf{defined} in procedures.}
	\end{frame}
	
	% ================================================================
	% SLIDE 75-78: GOVERNANCE ESCALATION AND OFFBOARDING
	% ================================================================
	\begin{frame}
		\frametitle{\fontsize{11}{13}\selectfont Q73: Governance Escalation}
		\begin{block}{\fontsize{8}{10}\selectfont Topic: Governance Escalation}
			\scriptsize
			How should investigations be escalated through governance?
		\end{block}
		\vspace{0.05cm}
		\stepbox{\scriptsize
			\keypoint{Correct Answer:} Through defined escalation paths to committees, senior management, and the board as appropriate.
		}
		\vspace{0.03cm}
		\tipbox{\scriptsize \textbf{Key Insight:} Escalation paths should be \textbf{clear} and \textbf{documented}.}
	\end{frame}
	
	\begin{frame}
		\frametitle{\fontsize{11}{13}\selectfont Q74: Decision-Making on Offboarding}
		\begin{block}{\fontsize{8}{10}\selectfont Topic: Offboarding Decisions}
			\scriptsize
			What should be considered when offboarding a customer?
		\end{block}
		\vspace{0.05cm}
		\stepbox{\scriptsize
			\keypoint{Correct Answer:} Risk level, regulatory requirements, potential for de-risking, and liaison with law enforcement/FIUs.
		}
		\vspace{0.03cm}
		\tipbox{\scriptsize \textbf{Key Insight:} Offboarding decisions should be \textbf{governed} and \textbf{documented}.}
	\end{frame}
	
	\begin{frame}
		\frametitle{\fontsize{11}{13}\selectfont Q75: Offboarding Documentation}
		\begin{block}{\fontsize{8}{10}\selectfont Topic: Offboarding Documentation}
			\scriptsize
			Why is offboarding documentation important?
		\end{block}
		\vspace{0.05cm}
		\stepbox{\scriptsize
			\keypoint{Correct Answer:} To demonstrate that offboarding was based on risk and not discriminatory.
		}
		\vspace{0.03cm}
		\tipbox{\scriptsize \textbf{Key Insight:} Documentation protects against \textbf{regulatory scrutiny}.}
	\end{frame}
	
	\begin{frame}
		\frametitle{\fontsize{11}{13}\selectfont Q76: FIU Liaison on Offboarding}
		\begin{block}{\fontsize{8}{10}\selectfont Topic: FIU Liaison}
			\scriptsize
			When should institutions liaise with FIUs on offboarding?
		\end{block}
		\vspace{0.05cm}
		\stepbox{\scriptsize
			\keypoint{Correct Answer:} When there are ongoing investigations or law enforcement interest in the customer.
		}
		\vspace{0.03cm}
		\tipbox{\scriptsize \textbf{Key Insight:} Coordinating with FIUs prevents \textbf{interference} with investigations.}
	\end{frame}
	
	% ================================================================
	% SLIDE 79-82: RESPONDING TO COURT ORDERS AND LAW ENFORCEMENT
	% ================================================================
	\begin{frame}
		\frametitle{\fontsize{11}{13}\selectfont Q77: Responding to Subpoenas}
		\begin{block}{\fontsize{8}{10}\selectfont Topic: Subpoenas}
			\scriptsize
			How should institutions respond to subpoenas?
		\end{block}
		\vspace{0.05cm}
		\stepbox{\scriptsize
			\keypoint{Correct Answer:} Respond promptly within legal timelines, document the response, and ensure confidentiality.
		}
		\vspace{0.03cm}
		\tipbox{\scriptsize \textbf{Key Insight:} Subpoenas are \textbf{legally binding} requests.}
	\end{frame}
	
	\begin{frame}
		\frametitle{\fontsize{11}{13}\selectfont Q78: Law Enforcement Requests}
		\begin{block}{\fontsize{8}{10}\selectfont Topic: Law Enforcement Requests}
			\scriptsize
			How should law enforcement requests be handled?
		\end{block}
		\vspace{0.05cm}
		\stepbox{\scriptsize
			\keypoint{Correct Answer:} Verify the request, document all communications, respond appropriately, and consult legal counsel if needed.
		}
		\vspace{0.03cm}
		\tipbox{\scriptsize \textbf{Key Insight:} \textbf{Documentation} of all interactions is essential.}
	\end{frame}
	
	\begin{frame}
		\frametitle{\fontsize{11}{13}\selectfont Q79: Section 314(a) Requests (US)}
		\begin{block}{\fontsize{8}{10}\selectfont Topic: Section 314(a)}
			\scriptsize
			What is a Section 314(a) request?
		\end{block}
		\vspace{0.05cm}
		\stepbox{\scriptsize
			\keypoint{Correct Answer:} A FinCEN request to financial institutions to search records for information about individuals or entities involved in a money laundering investigation.
		}
		\vspace{0.03cm}
		\tipbox{\scriptsize \textbf{Key Insight:} Response to 314(a) requests is \textbf{mandatory}.}
	\end{frame}
	
	\begin{frame}
		\frametitle{\fontsize{11}{13}\selectfont Q80: Section 314(b) Sharing}
		\begin{block}{\fontsize{8}{10}\selectfont Topic: Section 314(b)}
			\scriptsize
			What is Section 314(b) information sharing?
		\end{block}
		\vspace{0.05cm}
		\stepbox{\scriptsize
			\keypoint{Correct Answer:} Financial institutions sharing information with each other for AML purposes, subject to opt-in with FinCEN.
		}
		\vspace{0.03cm}
		\wrongbox{\scriptsize \textbf{Common Trap:} "Sharing is automatic" — INCORRECT. Institutions must \textbf{opt-in} before sharing.}
		\vspace{0.03cm}
		\tipbox{\scriptsize \textbf{Key Insight:} Opt-in is required \textbf{before} sharing.}
	\end{frame}
	
	% ================================================================
	% SLIDE 83-86: DE-RISKING AND FINANCIAL INCLUSION
	% ================================================================
	\begin{frame}
		\frametitle{\fontsize{11}{13}\selectfont Q81: De-risking Definition}
		\begin{block}{\fontsize{8}{10}\selectfont Topic: De-risking}
			\scriptsize
			What is de-risking in the AML context?
		\end{block}
		\vspace{0.05cm}
		\stepbox{\scriptsize
			\keypoint{Correct Answer:} The practice of terminating or restricting business relationships with entire categories of customers to avoid AML risk.
		}
		\vspace{0.03cm}
		\wrongbox{\scriptsize \textbf{Common Trap:} "De-risking is recommended by regulators" — INCORRECT. Regulators \textbf{discourage} indiscriminate de-risking.}
		\vspace{0.03cm}
		\tipbox{\scriptsize \textbf{Key Insight:} FATF has issued guidance \textbf{discouraging} de-risking.}
	\end{frame}
	
	\begin{frame}
		\frametitle{\fontsize{11}{13}\selectfont Q82: De-risking Consequences}
		\begin{block}{\fontsize{8}{10}\selectfont Topic: De-risking Consequences}
			\scriptsize
			What are consequences of de-risking?
		\end{block}
		\vspace{0.05cm}
		\stepbox{\scriptsize
			\keypoint{Correct Answer:} Financial exclusion, push to unregulated channels, and reputational damage.
		}
		\vspace{0.03cm}
		\tipbox{\scriptsize \textbf{Key Insight:} De-risking can \textbf{undermine} financial inclusion.}
	\end{frame}
	
	\begin{frame}
		\frametitle{\fontsize{11}{13}\selectfont Q83: Financial Inclusion and AML}
		\begin{block}{\fontsize{8}{10}\selectfont Topic: Financial Inclusion}
			\scriptsize
			How can financial inclusion be balanced with AML requirements?
		\end{block}
		\vspace{0.05cm}
		\stepbox{\scriptsize
			\keypoint{Correct Answer:} By applying proportionate controls and using simplified due diligence for lower-risk customers.
		}
		\vspace{0.03cm}
		\tipbox{\scriptsize \textbf{Key Insight:} Proportionality enables \textbf{both} inclusion and compliance.}
	\end{frame}
	
	\begin{frame}
		\frametitle{\fontsize{11}{13}\selectfont Q84: Alternative Verification Methods}
		\begin{block}{\fontsize{8}{10}\selectfont Topic: Alternative Verification}
			\scriptsize
			What alternative verification methods support inclusion?
		\end{block}
		\vspace{0.05cm}
		\stepbox{\scriptsize
			\keypoint{Correct Answer:} Biometrics, digital identity systems, and community-based verification.
		}
		\vspace{0.03cm}
		\tipbox{\scriptsize \textbf{Key Insight:} Technology can bridge the inclusion gap.}
	\end{frame}
	
	% ================================================================
	% SLIDE 87-90: CUSTOMER COMMUNICATION AND TIPPING-OFF
	% ================================================================
	\begin{frame}
		\frametitle{\fontsize{11}{13}\selectfont Q85: Customer Communication}
		\begin{block}{\fontsize{8}{10}\selectfont Topic: Customer Communication}
			\scriptsize
			How should institutions communicate with customers about AML requirements?
		\end{block}
		\vspace{0.05cm}
		\stepbox{\scriptsize
			\keypoint{Correct Answer:} Through established channels, clear explanations of requirements, and RFIs for additional information.
		}
		\vspace{0.03cm}
		\tipbox{\scriptsize \textbf{Key Insight:} Communication should be \textbf{professional} and \textbf{respectful}.}
	\end{frame}
	
	\begin{frame}
		\frametitle{\fontsize{11}{13}\selectfont Q86: Requests for Information (RFIs)}
		\begin{block}{\fontsize{8}{10}\selectfont Topic: RFIs}
			\scriptsize
			What is an RFI in the AML context?
		\end{block}
		\vspace{0.05cm}
		\stepbox{\scriptsize
			\keypoint{Correct Answer:} A request to a customer for additional information to satisfy CDD/EDD requirements.
		}
		\vspace{0.03cm}
		\tipbox{\scriptsize \textbf{Key Insight:} RFIs should be \textbf{clear} and \textbf{justified}.}
	\end{frame}
	
	\begin{frame}
		\frametitle{\fontsize{11}{13}\selectfont Q87: Tipping-Off Consequences}
		\begin{block}{\fontsize{8}{10}\selectfont Topic: Tipping-Off}
			\scriptsize
			What is tipping-off and what are its consequences?
		\end{block}
		\vspace{0.05cm}
		\stepbox{\scriptsize
			\keypoint{Correct Answer:} Tipping-off is informing a subject that they are under investigation. It is a criminal offense in many jurisdictions.
		}
		\vspace{0.03cm}
		\tipbox{\scriptsize \textbf{Key Insight:} Tipping-off can \textbf{destroy} an investigation.}
	\end{frame}
	
	\begin{frame}
		\frametitle{\fontsize{11}{13}\selectfont Q88: Raising Customer Awareness}
		\begin{block}{\fontsize{8}{10}\selectfont Topic: Raising Awareness}
			\scriptsize
			How can institutions raise awareness of AML obligations among customers?
		\end{block}
		\vspace{0.05cm}
		\stepbox{\scriptsize
			\keypoint{Correct Answer:} Through onboarding communications, periodic notices, and online resources.
		}
		\vspace{0.03cm}
		\tipbox{\scriptsize \textbf{Key Insight:} Educated customers are \textbf{more cooperative}.}
	\end{frame}
	
	% ================================================================
	% SLIDE 91-94: TRAINING AND AWARENESS
	% ================================================================
	\begin{frame}
		\frametitle{\fontsize{11}{13}\selectfont Q89: AML Training Requirements}
		\begin{block}{\fontsize{8}{10}\selectfont Topic: AML Training}
			\scriptsize
			What are the AML training requirements for staff?
		\end{block}
		\vspace{0.05cm}
		\stepbox{\scriptsize
			\keypoint{Correct Answer:} Role-based training, regular refresher training, and updates on regulatory changes.
		}
		\vspace{0.03cm}
		\tipbox{\scriptsize \textbf{Key Insight:} Training should be \textbf{continuous}, not a one-time event.}
	\end{frame}
	
	\begin{frame}
		\frametitle{\fontsize{11}{13}\selectfont Q90: Role-Based Training}
		\begin{block}{\fontsize{8}{10}\selectfont Topic: Role-Based Training}
			\scriptsize
			Why is role-based training important?
		\end{block}
		\vspace{0.05cm}
		\stepbox{\scriptsize
			\keypoint{Correct Answer:} Different employees need different knowledge based on their roles and exposure to risks.
		}
		\vspace{0.03cm}
		\tipbox{\scriptsize \textbf{Key Insight:} Frontline staff need different training than investigators.}
	\end{frame}
	
	\begin{frame}
		\frametitle{\fontsize{11}{13}\selectfont Q91: Training Effectiveness}
		\begin{block}{\fontsize{8}{10}\selectfont Topic: Training Effectiveness}
			\scriptsize
			How is training effectiveness measured?
		\end{block}
		\vspace{0.05cm}
		\stepbox{\scriptsize
			\keypoint{Correct Answer:} Through assessments, employee feedback, and observation of application in the workplace.
		}
		\vspace{0.03cm}
		\tipbox{\scriptsize \textbf{Key Insight:} Effectiveness is measured by \textbf{application}, not just completion.}
	\end{frame}
	
	\begin{frame}
		\frametitle{\fontsize{11}{13}\selectfont Q92: Raising Staff Awareness}
		\begin{block}{\fontsize{8}{10}\selectfont Topic: Staff Awareness}
			\scriptsize
			How can staff awareness of AML risks be raised?
		\end{block}
		\vspace{0.05cm}
		\stepbox{\scriptsize
			\keypoint{Correct Answer:} Through regular training, communications, and sharing of typologies and case studies.
		}
		\vspace{0.03cm}
		\tipbox{\scriptsize \textbf{Key Insight:} Awareness is the \textbf{foundation} of detection.}
	\end{frame}
	
	% ================================================================
	% SLIDE 95-98: AFC TEAM STRUCTURE
	% ================================================================
	\begin{frame}
		\frametitle{\fontsize{11}{13}\selectfont Q93: AFC Team Structure}
		\begin{block}{\fontsize{8}{10}\selectfont Topic: AFC Team Structure}
			\scriptsize
			What is the typical structure of an AFC team?
		\end{block}
		\vspace{0.05cm}
		\stepbox{\scriptsize
			\keypoint{Correct Answer:} MLRO/BSA Officer, compliance officers, investigators, analysts, and support staff.
		}
		\vspace{0.03cm}
		\tipbox{\scriptsize \textbf{Key Insight:} Structure depends on the \textbf{institution's size} and \textbf{complexity}.}
	\end{frame}
	
	\begin{frame}
		\frametitle{\fontsize{11}{13}\selectfont Q94: MLRO/BSA Officer Role}
		\begin{block}{\fontsize{8}{10}\selectfont Topic: MLRO/BSA Officer}
			\scriptsize
			What is the role of the MLRO/BSA Officer?
		\end{block}
		\vspace{0.05cm}
		\stepbox{\scriptsize
			\keypoint{Correct Answer:} Designated individual responsible for oversight of the AML program, filing SARs, and acting as the point of contact for regulators.
		}
		\vspace{0.03cm}
		\tipbox{\scriptsize \textbf{Key Insight:} The MLRO must have \textbf{access to the board}.}
	\end{frame}
	
	\begin{frame}
		\frametitle{\fontsize{11}{13}\selectfont Q95: AFC Investigator Role}
		\begin{block}{\fontsize{8}{10}\selectfont Topic: Investigator Role}
			\scriptsize
			What is the role of an AFC investigator?
		\end{block}
		\vspace{0.05cm}
		\stepbox{\scriptsize
			\keypoint{Correct Answer:} Investigating alerts, gathering information, documenting findings, and recommending SAR filing.
		}
		\vspace{0.03cm}
		\tipbox{\scriptsize \textbf{Key Insight:} Investigators are the \textbf{core} of the AFC function.}
	\end{frame}
	
	\begin{frame}
		\frametitle{\fontsize{11}{13}\selectfont Q96: Team Competencies}
		\begin{block}{\fontsize{8}{10}\selectfont Topic: Team Competencies}
			\scriptsize
			What competencies are needed in an AFC team?
		\end{block}
		\vspace{0.05cm}
		\stepbox{\scriptsize
			\keypoint{Correct Answer:} Regulatory knowledge, analytical skills, investigation experience, and communication skills.
		}
		\vspace{0.03cm}
		\tipbox{\scriptsize \textbf{Key Insight:} \textbf{Diverse skills} are needed for effective team performance.}
	\end{frame}
	
	% ================================================================
	% SLIDE 99-102: AFC AND FRONT OFFICE
	% ================================================================
	\begin{frame}
		\frametitle{\fontsize{11}{13}\selectfont Q97: AFC and Front Office Interaction}
		\begin{block}{\fontsize{8}{10}\selectfont Topic: AFC and Front Office}
			\scriptsize
			How should AFC professionals interact with the front office?
		\end{block}
		\vspace{0.05cm}
		\stepbox{\scriptsize
			\keypoint{Correct Answer:} Collaboratively, with clear roles, escalation paths, and mutual respect for each other's responsibilities.
		}
		\vspace{0.03cm}
		\tipbox{\scriptsize \textbf{Key Insight:} \textbf{Collaboration} is essential for effective compliance.}
	\end{frame}
	
	\begin{frame}
		\frametitle{\fontsize{11}{13}\selectfont Q98: Relationship Manager Role}
		\begin{block}{\fontsize{8}{10}\selectfont Topic: Relationship Managers}
			\scriptsize
			What is the role of relationship managers in AFC?
		\end{block}
		\vspace{0.05cm}
		\stepbox{\scriptsize
			\keypoint{Correct Answer:} Identify suspicious activity, escalate concerns, and support the AFC team in investigations.
		}
		\vspace{0.03cm}
		\tipbox{\scriptsize \textbf{Key Insight:} Relationship managers are the \textbf{eyes and ears} of AFC.}
	\end{frame}
	
	\begin{frame}
		\frametitle{\fontsize{11}{13}\selectfont Q99: Escalation from Front Office}
		\begin{block}{\fontsize{8}{10}\selectfont Topic: Front Office Escalation}
			\scriptsize
			When should the front office escalate to AFC?
		\end{block}
		\vspace{0.05cm}
		\stepbox{\scriptsize
			\keypoint{Correct Answer:} When they observe unusual activity, have concerns about a customer, or need guidance on AML requirements.
		}
		\vspace{0.03cm}
		\tipbox{\scriptsize \textbf{Key Insight:} Escalation should be \textbf{encouraged}, not discouraged.}
	\end{frame}
	
	\begin{frame}
		\frametitle{\fontsize{11}{13}\selectfont Q100: Front Office Training}
		\begin{block}{\fontsize{8}{10}\selectfont Topic: Front Office Training}
			\scriptsize
			What training should be provided to the front office?
		\end{block}
		\vspace{0.05cm}
		\stepbox{\scriptsize
			\keypoint{Correct Answer:} Identification of red flags, escalation procedures, and confidentiality obligations.
		}
		\vspace{0.03cm}
		\tipbox{\scriptsize \textbf{Key Insight:} Front office staff need \textbf{practical} training.}
	\end{frame}
	
	% ================================================================
	% SLIDE 103-106: OTHER SUPPORTING FUNCTIONS
	% ================================================================
	\begin{frame}
		\frametitle{\fontsize{11}{13}\selectfont Q101: Data Privacy and AFC}
		\begin{block}{\fontsize{8}{10}\selectfont Topic: Data Privacy}
			\scriptsize
			How does data privacy support the AFC program?
		\end{block}
		\vspace{0.05cm}
		\stepbox{\scriptsize
			\keypoint{Correct Answer:} Ensures data is collected and processed in compliance with privacy laws, protecting the institution from legal risk.
		}
		\vspace{0.03cm}
		\tipbox{\scriptsize \textbf{Key Insight:} Data privacy and AFC are \textbf{interconnected}.}
	\end{frame}
	
	\begin{frame}
		\frametitle{\fontsize{11}{13}\selectfont Q102: Cybersecurity and AFC}
		\begin{block}{\fontsize{8}{10}\selectfont Topic: Cybersecurity}
			\scriptsize
			How does cybersecurity support the AFC program?
		\end{block}
		\vspace{0.05cm}
		\stepbox{\scriptsize
			\keypoint{Correct Answer:} Protects AFC data and systems from cyber threats and provides intelligence on cyber-enabled financial crime.
		}
		\vspace{0.03cm}
		\tipbox{\scriptsize \textbf{Key Insight:} Cybersecurity and AFC have \textbf{overlapping} risks.}
	\end{frame}
	
	\begin{frame}
		\frametitle{\fontsize{11}{13}\selectfont Q103: Risk Management Functions}
		\begin{block}{\fontsize{8}{10}\selectfont Topic: Risk Management}
			\scriptsize
			How do other risk management functions support AFC?
		\end{block}
		\vspace{0.05cm}
		\stepbox{\scriptsize
			\keypoint{Correct Answer:} Through integrated risk assessments, shared intelligence, and coordinated responses.
		}
		\vspace{0.03cm}
		\tipbox{\scriptsize \textbf{Key Insight:} Risk management functions should be \textbf{integrated}.}
	\end{frame}
	
	\begin{frame}
		\frametitle{\fontsize{11}{13}\selectfont Q104: Integrated AFC Program}
		\begin{block}{\fontsize{8}{10}\selectfont Topic: Integrated Program}
			\scriptsize
			What makes an integrated AFC program robust?
		\end{block}
		\vspace{0.05cm}
		\stepbox{\scriptsize
			\keypoint{Correct Answer:} Collaboration across functions, shared risk intelligence, and coordinated controls.
		}
		\vspace{0.03cm}
		\tipbox{\scriptsize \textbf{Key Insight:} Integration \textbf{strengthens} the overall program.}
	\end{frame}
	
	% ================================================================
	% SLIDE 107-110: SUMMARY
	% ================================================================
	\begin{frame}
		\frametitle{\fontsize{11}{13}\selectfont SUMMARY: Key Principles — Domain C}
		\scriptsize
		\begin{block}{\fontsize{9}{11}\selectfont Building an Anti-Financial Crime Compliance Program:}
			\vspace{0.05cm}
			\scriptsize
			\begin{tabular}{|p{0.35\textwidth}|p{0.60\textwidth}|}
				\hline
				\keypoint{Concept} & \keypoint{Application} \\
				\hline
				Five Pillars & Internal controls, compliance officer, training, testing, CDD \\
				\hline
				Three Lines of Defense & First: Operations; Second: Compliance; Third: Audit \\
				\hline
				Risk Appetite Statement & Defines risk tolerance; guides decisions \\
				\hline
				Enterprise Risk Assessment & Inherent risk; controls; residual risk \\
				\hline
				Risk-Based Approach & Proportional controls; resource allocation \\
				\hline
				Policies/Standards/Procedures & What; must; how \\
				\hline
				Horizon Scanning & Identifying emerging risks \\
				\hline
				Governing Committees & Oversight; decision-making \\
				\hline
				KPIs/KRIs & Performance; risk indicators \\
				\hline
				Reporting Requirements & SARs; CTRs; jurisdictional differences \\
				\hline
				Operational/Financial/Non-Financial Risk & AFC failure impacts all \\
				\hline
				Onboarding/Lifecycle Controls & CDD; EDD; monitoring; exit controls \\
				\hline
				High-Risk Scenarios & Conflict zones; dual-use; PEPs; complex structures \\
				\hline
				Other Due Diligence & Employees; vendors; third parties \\
				\hline
				Investigations & Documentation; SAR writing; escalation \\
				\hline
				Offboarding & Governance; FIU liaison; documentation \\
				\hline
				Court Orders/Law Enforcement & Subpoenas; 314(a); 314(b) \\
				\hline
				De-risking vs Inclusion & Discouraged; use proportionate controls \\
				\hline
				Tipping-Off & Criminal offense; don't alert subjects \\
				\hline
				Training & Role-based; continuous; effective \\
				\hline
				AFC Team & MLRO; investigators; analysts \\
				\hline
				Front Office & Eyes and ears; escalation \\
				\hline
				Supporting Functions & Data privacy; cybersecurity; risk management \\
				\hline
			\end{tabular}
			\vspace{0.05cm}
		\end{block}
		\vspace{0.05cm}
		\tipbox{\scriptsize \textbf{FINAL LESSON:} Domain C tests understanding of \textbf{how to build and run} an AFC program — from governance to operations to reporting. Understand the \textit{structure} and the \textit{answers} become obvious.}
	\end{frame}
	
	\begin{frame}
		\frametitle{\fontsize{11}{13}\selectfont EXAM TIPS — Domain C}
		\scriptsize
		\begin{block}{\fontsize{9}{11}\selectfont Common Traps to Avoid:}
			\vspace{0.05cm}
			\scriptsize
			\begin{itemize}
				\item \highlight{Five Pillars}: Remember the \textbf{5th pillar} (CDD) added in 2016.
				\item \highlight{Three Lines}: First = \textbf{Operations}; Second = \textbf{Compliance}; Third = \textbf{Audit}.
				\item \highlight{Audit Independence}: Audit \textbf{cannot} design controls it tests.
				\item \highlight{Risk-Based Approach}: Not less compliance — \textbf{targeted} compliance.
				\item \highlight{De-risking}: \textbf{Discouraged} by regulators.
				\item \highlight{Tipping-Off}: \textbf{Criminal offense} in many jurisdictions.
				\item \highlight{314(b)}: \textbf{Opt-in required} before sharing.
				\item \highlight{Policies vs Procedures}: Policies = "what"; Procedures = "how".
				\item \highlight{MLRO/BSA Officer}: Must have \textbf{board access}.
				\item \highlight{EDD for PEPs}: \textbf{Mandatory}, not optional.
			\end{itemize}
			\vspace{0.05cm}
		\end{block}
		\vspace{0.05cm}
		\tipbox{\scriptsize \textbf{Key Insight:} Focus on understanding the \textit{structure} and \textit{governance} of an AFC program.}
	\end{frame}
	
	\begin{frame}
		\centering
		\vspace{2.5cm}
		{\fontsize{16}{19}\selectfont \textbf{Thank You}} \\[0.2cm]
		{\fontsize{11}{13}\selectfont Keep learning. Keep growing.} \\[0.15cm]
		\rule{3cm}{0.5pt} \\[0.15cm]
		{\fontsize{8}{10}\selectfont \textit{Building an effective AFC program requires governance, controls, and continuous improvement.}}
		\vspace{0.3cm}
		{\fontsize{7}{9}\selectfont \textcolor{darkgray}{End of Domain C Review — 110 Slides}}
	\end{frame}
	
\end{document}